\documentclass{article}

\usepackage{booktabs}
\usepackage{tabularx}
\usepackage{hyperref}
\usepackage{graphicx}

\hypersetup{
    colorlinks=true,       % false: boxed links; true: colored links
    linkcolor=red,          % color of internal links (change box color with linkbordercolor)
    citecolor=green,        % color of links to bibliography
    filecolor=magenta,      % color of file links
    urlcolor=cyan           % color of external links
}

\title{Hazard Analysis\\\progname}

\author{\authname}

\date{}

\input{../Comments.text}
\input{../Common.text}

\begin{document}

\maketitle
\thispagestyle{empty}

~\newpage

\pagenumbering{roman}

\begin{table}[hp]
\caption{Revision History} \label{TblRevisionHistory}
\begin{tabularx}{\textwidth}{llX}
\toprule
\textbf{Date} & \textbf{Developer(s)} & \textbf{Change}\\
\midrule
Date1 & Name(s) & Description of changes\\
Date2 & Name(s) & Description of changes\\
... & ... & ...\\
\bottomrule
\end{tabularx}
\end{table}

~\newpage

\tableofcontents

~\newpage

\pagenumbering{arabic}

\wss{You are free to modify this template.}

\section{Introduction}

\wss{You can include your definition of what a hazard is here.}

\subsection{Background}

The Troitsky Toolbox is a web-based application that supports the McMaster
Troitsky club, a student organization whose teams design, analyze, and build
bridges out of popsicle sticks, white glue, and dental floss for the annual
Troitsky Bridge Building Competition. The Toolbox centralizes the club's
knowledge, including past designs, material testing data, and exit documents,
and provides technical tools that allow teams to visualize their bridge models,
estimate their performance, and check them against the competition rules. The
problem and goals of the project are described in more detail in the Problem
Statement and Goals document, and the system itself is specified in the
Software Requirements Specification (SRS).

\subsection{Hazard Definition}

Following Leveson, a hazard is defined as a system state or set of conditions
that, together with a particular set of worst-case environmental conditions,
will lead to a loss. For the Troitsky Toolbox, a hazard is therefore any
condition of the system that, combined with how members use it, could lead to
the loss of club knowledge, the disclosure or corruption of club data, the
disruption of club activities, or members acting on incorrect information.

The Toolbox does not control any physical equipment, so the hazards it presents
are primarily to the club's data and decisions rather than to people. The one
exception is indirect: the Performance Analysis System produces estimates of a
bridge's failure load, and a team that trusts an incorrect estimate may build or
load test a bridge on that basis. This document identifies these hazards, the
failure modes that cause them, and the safety and security requirements that
mitigate them.

\section{Scope and Purpose of Hazard Analysis}

\wss{You should say what \textbf{loss} could be incurred because of the
hazards.}

The purpose of this hazard analysis is to identify the ways in which the
Troitsky Toolbox could fail or be misused, assess their effects on the club and
its members, and derive requirements that prevent or mitigate them. The
resulting safety and security requirements are added to the SRS, so that they
are designed, implemented, and verified alongside the rest of the system.

The losses that could result from the hazards identified in this document are:

\begin{itemize}
    \item \textbf{Loss of club knowledge:} Designs, test data, or exit documents
    are deleted, corrupted, or never captured, so that the knowledge the Toolbox
    exists to preserve is lost when members graduate. As the Toolbox becomes the
    club's central repository, a single failure can affect several years of
    accumulated information.
    \item \textbf{Loss of data integrity:} Stored data, such as material test
    results or competition rules, is incorrect or inconsistent, so that the
    results of the analysis tools are wrong.
    \item \textbf{Misleading analysis results:} A team relies on an incorrect
    performance estimate or compliance result, leading to a weaker bridge, a
    bridge that is disqualified for violating a competition rule, or wasted
    build time and materials.
    \item \textbf{Loss of confidentiality:} Members' personal information or the
    club's unpublished designs are disclosed to unauthorized users, including
    other competing universities.
    \item \textbf{Loss of availability:} The Toolbox, or one of its tools, is
    unavailable when members need it, such as in the weeks leading up to the
    competition.
    \item \textbf{Loss of trust:} Repeated failures lead members to stop using
    the Toolbox and return to scattered, undocumented information sharing,
    defeating the purpose of the project.
\end{itemize}

The analysis covers every software component of the Toolbox listed in
Section~\ref{SecBoundaries}, and the interfaces between the Toolbox and the
external components it depends on. The following are outside the scope of this
analysis:

\begin{itemize}
    \item The physical construction and load testing of bridges, including the
    safety procedures followed by the club during testing.
    \item The CAD software that teams use to create their bridge models.
    \item The Troitsky Bridge Building Competition itself, including its rules
    and judging.
    \item Failures of the hosting infrastructure's physical hardware, beyond
    their effect on the availability and durability of the Toolbox's data.
\end{itemize}

\section{System Boundaries and Components} \label{SecBoundaries}

\wss{Dividing the system into components will help you brainstorm the hazards.
You shouldn't do a full design of the components, just get a feel for the major
ones.  For projects that involve hardware, the components will typically include
each individual piece of hardware.  If your software will have a database, or an
important library, these are also potential components.}

The Troitsky Toolbox is structured around a central broker: each tool, as well
as the shared data storage, is a service that registers with the broker, and all
communication between services passes through it. The components below are
those described in Section S.1 of the SRS, and are shown in
Figure~\ref{FigComponents}. The requirements each component is responsible for
are listed in Section S.2 of the SRS, using the prefixes given below.

\begin{figure}[htbp]
    \centering
    \includegraphics[width=\textwidth]{../SRS/System/s1component-dark.png}
    \caption{Components of the Troitsky Toolbox and the interfaces through which
    they interact.}
    \label{FigComponents}
\end{figure}

\subsection{Inside the System Boundary}

The following components are developed by the team and are within the system's
control and responsibility:

\begin{itemize}
    \item \textbf{Web Server (WS):} Serves the Toolbox Web App and the public
    landing page, handles member login, and forwards requests for club services
    to the Toolbox Broker.
    \item \textbf{Access Control System (AC):} Manages member accounts and
    roles, authenticates members, and decides which services and data each role
    may access.
    \item \textbf{Toolbox Broker (BR):} Registers services, verifies every
    request with the Access Control System, and routes it to the service that
    provides the requested operation. Because all communication passes through
    the broker, a failure of the broker affects every tool.
    \item \textbf{Knowledge Repository (KR):} Stores all club information,
    including documents, bridge models, material test results, competition
    rules, and handoff documents, and allows members to search it.
    \item \textbf{Model Analysis System (MA):} Imports bridge models in the STEP
    format, renders them in 3D, and extracts the bridge properties needed for
    analysis, such as its dimensions and mass.
    \begin{itemize}
        \item \textbf{Performance Analysis System (PA):} Estimates the failure
        load and load-to-mass ratio of a bridge model, using the extracted model
        data and the material properties provided by the Material Testing
        Manager.
        \item \textbf{Rule Compliance System (RC):} Checks a bridge model against
        the competition rules configured by club members.
    \end{itemize}
    \item \textbf{Material Testing Manager (MT):} Imports and stores material
    test results from CSV and Excel files, compares them against previous years,
    and derives the material properties used by the Performance Analysis System.
    \item \textbf{Knowledge Transfer Manager (KT):} Collects exit documents and
    role records from executives and team leads, and stores them in the
    Knowledge Repository.
\end{itemize}

\subsection{Outside the System Boundary}

The following entities interact with the system, but are not under its direct
control:

\begin{itemize}
    \item \textbf{Club members:} The general members, executives, and software
    subteam maintainers who use the Toolbox. Their decisions, including how they
    act on the results of the analysis tools, are outside the system's control.
    \item \textbf{CAD software and STEP files:} The software teams use to design
    their bridges, and the STEP files it exports. The Toolbox receives these
    files but does not control how they are created.
    \item \textbf{CAD/geometry library:} The external library used by the Model
    Analysis System to read STEP files.
    \item \textbf{Material test data files:} The CSV and Excel files produced by
    members during material testing, whose content depends on how the tests
    were performed and recorded.
    \item \textbf{Competition rules:} The official rules published by the
    competition organizers, which change from year to year and must be
    configured into the Rule Compliance System by members.
    \item \textbf{Hosting infrastructure:} The server, operating system, database
    server, and network on which the Toolbox is deployed.
    \item \textbf{Web browsers:} The browsers through which members access the
    Toolbox Web App.
\end{itemize}

\section{Critical Assumptions}

\wss{These assumptions that are made about the software or system.  You should
minimize the number of assumptions that remove potential hazards.  For instance,
you could assume a part will never fail, but it is generally better to include
this potential failure mode.}

The following assumptions are made about the system and its environment. Each
assumption states the class of hazards it removes from the scope of this
analysis. Assumptions have been kept to a minimum, and failures that the system
can reasonably detect or mitigate, such as invalid input files, service
outages, and data entry errors, remain within the scope of the analysis.

\begin{enumerate}
    \item \textbf{Analysis results are advisory:} Members treat the results of
    the Performance Analysis and Rule Compliance Systems as guidance, and the
    club verifies the final design against the official competition rules and
    follows its own safety procedures when load testing bridges. This assumption
    removes physical injury during bridge testing from the scope of this
    analysis. Hazards caused by incorrect or misleading results, and by results
    being presented in a way that encourages over-reliance, remain in scope.

    \item \textbf{Members act in good faith:} Authenticated members do not
    intentionally misuse their access, for example by deliberately deleting or
    falsifying club data. This assumption removes deliberate insider sabotage
    from the scope of this analysis. Accidental misuse by members, and attacks
    by people outside the club, remain in scope.

    \item \textbf{STEP files reflect the intended design:} The STEP files that
    teams upload accurately represent the bridge they intend to build. This
    assumption removes differences between a team's CAD model and its physical
    bridge from the scope of this analysis. Files that are invalid, corrupted, or
    that cannot be read remain in scope.

    \item \textbf{The CAD/geometry library is correct:} The external library used
    by the Model Analysis System correctly reads the geometry of valid STEP
    files. This assumption removes defects within the library itself from the
    scope of this analysis. Errors in how the Toolbox uses the library's output
    remain in scope.

    \item \textbf{Test data reflects the tests performed:} The material test
    results that members import were measured using the club's testing
    procedure. This assumption removes errors in the physical testing process
    from the scope of this analysis. Data entry errors, malformed files, and
    physically impossible values remain in scope.

    \item \textbf{Hosting infrastructure is maintained:} The server, operating
    system, and database server on which the Toolbox is deployed are kept
    updated and secured by whoever hosts them. This assumption removes
    vulnerabilities within the hosting platform itself from the scope of this
    analysis. Outages of the hosting infrastructure, and their effect on the
    availability and durability of the Toolbox's data, remain in scope.
\end{enumerate}

\section{Failure Mode and Effect Analysis}

\wss{Include your FMEA table here. This is the most important part of this document.}
\wss{The safety requirements in the table do not have to have the prefix SR.
The most important thing is to show traceability to your SRS. You might trace to
requirements you have already written, or you might need to add new
requirements.}
\wss{If no safety requirement can be devised, other mitigation strategies can be
entered in the table, including strategies involving providing additional
documentation, and/or test cases.}

\section{Safety and Security Requirements}

\wss{Newly discovered requirements.  These should also be added to the SRS.  (A
rationale design process how and why to fake it.)}

\section{Roadmap}

\wss{Which safety requirements will be implemented as part of the capstone timeline?
Which requirements will be implemented in the future?}

\newpage{}

\section{Appendix --- Generative AI Use Disclosure}

\wss{ChatGPT (version 5) was used to brainstorm ideas for the template for this disclose section.}

\subsection{AI Tools Used}

\wss{Give the name of the AI tool(s) used and the version.}

\subsection{How and Where Used}

\wss{Explain what the tool(s) were used for.  Examples include: brainstorming,
explaining a technical concept, reviewing the document, generating or modifying
text, generating or modifying diagrams, identifying errors or omissions or
inconsistencies.}

\wss{For each of the uses, identify which parts of the document were affected.}

\wss{You do not need to report every interaction, but you should disclose the 
uses that materially contributed to your work}

\subsection{Verification of AI-Generated Content}

\wss{\begin{itemize}
    \item Describe how the team checked AI-generated or AI-assisted material
    before including it in the document.
    \item In particular, verify factual claims, technical assertions,
    requirements, calculations, references, and other information for which an
    AI system may produce incorrect or fabricated results.
    \item \textbf{The team is responsible for the accuracy, completeness,
consistency, and appropriateness of the submitted document, regardless of
whether material was generated by a human or an AI system.}
\end{itemize}}

\newpage{}

\section{Appendix --- Reflection}

\wss{Not required for CAS 741}

\input{../Reflection.text}

\begin{enumerate}
    \item What went well while writing this deliverable? 
    \item What pain points did you experience during this deliverable, and how
    did you resolve them?
    \item Which of your listed risks had your team thought of before this
    deliverable, and which did you think of while doing this deliverable? For
    the latter ones (ones you thought of while doing the Hazard Analysis), how
    did they come about?
    \item Other than the risk of physical harm (some projects may not have any
    appreciable risks of this form), list at least 2 other types of risk in
    software products. Why are they important to consider?
\end{enumerate}

\end{document}